% =====================================================================
%  1 Introduction
%  Paragraph plan P1 to P5 from reports/thesis_proposal.md, updated with
%  the final results. Literature notes: reports/key_points_thesis_proposal.md
%  and the Notion export in reports/.
% =====================================================================
\section{Introduction}
\label{sec:intro}

% ---- P1: the problem -------------------------------------------------
% - Offensiveness is subjective; single-label gold standards erase real
%   variation ("disagreement is signal, not noise", Denton et al. 2021;
%   Uma et al. 2021; Plank 2022).
% - D3CODE (Davani et al. 2024): raters from 21 countries, 8 cultural
%   regions, moral foundations (MFQ-2).
% - Demographics alone explain under 10% of rating variance (Hu & Collier
%   2024), so any group signal is small.
\fillin{P1: offensiveness is subjective; why single labels are a problem; D3CODE.}
\citep{denton2021whose,uma2021learning,plank2022problem,davani2024d3code,hu2024quantifying}

% ---- P2: LLMs in this role -------------------------------------------
% - LLMs are used as annotators / moderators; accurate on unanimous items,
%   overconfident on contested ones (Lu et al. 2025).
% - Persona prompting gives little, and more attributes can hurt
%   (Kamruzzaman et al. 2025; Giorgi 2025; Sarumi et al. 2025); LLMs fail
%   to predict which group finds something more harmful (Movva et al. 2024).
% - Own evidence: both models over-flag on all 4,519 D3CODE texts
%   (flag 79% / 85% vs 22% majority-offensive; recall 0.99, precision
%   0.27 / 0.25, AUC 0.83 / 0.82). Source: reports/v3_identity_vs_baseline.md 5.
%   (The proposal's "confidence 0.93 vs 0.77 on contested vs agreed
%   thirds" is from the old seeded sample: recompute on all texts or drop,
%   see CLAUDE.md section 10.)
\fillin{P2: LLMs as offensiveness judges and persona prompting; what is known; our over-flagging numbers.}

% ---- P3: the gap -----------------------------------------------------
% - Prior work is behavioural: does the output change with a persona?
%   It rarely asks how the cue is used inside the model.
% - It rarely includes a content-free control: any inserted sentence can
%   move a prediction (recency bias, Zhao et al. 2021; position effects,
%   Liu et al. 2024). Phrase length matters too (our v2 result: longer
%   placebos shift less).
% - Interpretability tools exist (attention analysis, value-weighted
%   attention, PASTA, activation steering) but have not been applied to
%   demographic framing in offensiveness judgments.
\fillin{P3: what prior work does not do: internal mechanism, content-free and length-matched controls.}

% ---- P4: what we do --------------------------------------------------
% - Two open models, the same 500 borderline D3CODE texts, 22 conditions
%   (14 demographic, 8 length-matched placebos) x 10 prompt formats, plus
%   8 place-name controls.
% - Behaviour: shift in log-odds against a length-adjusted placebo null.
% - Mechanism: per-head attention (raw and value-weighted), then causal
%   single-head profiling with PASTA against two nulls.
% - Steering on held-out texts: dampener / amplifier heads, steering
%   towards the text, activation steering, random-head controls, upper
%   bounds; success = closer to the real group's share (supervisor).
% - Finally: is there any group signal in D3CODE to align with?
\fillin{P4: our approach in one paragraph.}

% ---- RQs ---------------------------------------------------------------
% Current wording (CLAUDE.md section 2). Proposal wording for reference:
%   RQ1 Is the effect of a demographic framing distinguishable from the
%       effect of any inserted sentence of matched length?
%   RQ2 Where does the framing register (layers, heads), and does attention
%       to it predict the behavioural shift?
%   RQ3 Can steering toward the perspective tokens make the model's
%       judgments track the framed group's real ratings better than plain
%       prompting?
We address three research questions:
\begin{description}[leftmargin=1.2cm,style=nextline]
  \item[RQ1] Do LLMs change their offensiveness judgments when given a
    demographic identity, beyond what content-free phrases of the same
    length do?
  \item[RQ2] Do these changes match how real members of those groups judge
    the same texts?
  \item[RQ3] Where in the model does the identity effect arise, and can
    attention steering control it and bring the model closer to the real
    group ratings?
\end{description}

% ---- P5: contributions and outline -------------------------------------
% Candidate contributions (from the proposal, updated):
%  1. A controlled protocol for evaluating demographic prompting:
%     shared items for both models, a baseline band against ceiling
%     effects, exact token-length matching, a placebo pool with a
%     length-adjusted null, place-name controls, phrase as the
%     statistical unit.
%  2. Behavioural findings: only region carries content; half of it (Llama)
%     or all of it (Qwen, perspective formats) is the place name; framing
%     moves models away from real raters.
%  3. Mechanistic findings: a handful of middle-layer heads carry the
%     effect causally; it reaches the answer via intermediate tokens;
%     attention-pattern selection agrees poorly with causal selection.
%  4. A steering evaluation that separates alignment from calibration
%     (constant-shift baseline, ranking, flag rate).
%  5. A reliability analysis showing D3CODE has almost no item-level group
%     signal (demographic or moral), only small average offsets.
Our contributions are:
\begin{enumerate}
  \item \fillin{Protocol contribution.}
  \item \fillin{Behavioural finding.}
  \item \fillin{Mechanistic finding.}
  \item \fillin{Steering finding.}
  \item \fillin{Data / group-signal finding.}
\end{enumerate}

\fillin{Outline: \cref{sec:related} reviews \ldots, \cref{sec:method} describes \ldots,
\cref{sec:behaviour,sec:mechanism,sec:steering,sec:groupsignal} report \ldots,
\cref{sec:discussion} discusses \ldots}
